% The factorization-pattern setup: K over E over ℚ matched with
% {1} ⊆ H ⊆ G, and the primes 𝔓, 𝔭₁⋯𝔭_g, (p) with f₁+⋯+f_g = n.
% Source: book/tex/alg-NT/frobenius.tex (§ Frobenius elements control
% factorization) — tikzcd.
\documentclass[border=2mm]{standalone}
\usepackage{tikz}
\usetikzlibrary{arrows.meta}
\usepackage{amssymb}
\usepackage{amsmath}
\begin{document}
\begin{tikzpicture}[>=latex]
  \node (K)  at (0,3.0)   {$K$};
  \node (E)  at (0,1.5)   {$E$};
  \node (Q)  at (0,0)     {$\mathbb{Q}$};
  \node (e)  at (2.2,3.0) {$\{1\}$};
  \node (H)  at (2.2,1.5) {$H$};
  \node (G)  at (2.2,0)   {$G$};
  \node (PP) at (4.8,3.0) {$\operatorname{Frob}_{\mathfrak{P}}$};
  \node (pp) at (4.8,1.5) {$\mathfrak{p}_1 \dots \mathfrak{p}_g$};
  \node (p)  at (4.8,0)   {$(p)$};
  \node (f)  at (7.9,1.5) {$f_1 + \dots + f_g = n$};
  \draw (K) -- (E) node[midway,left] {};
  \draw (E) -- (Q) node[midway,left] {$n$};
  \draw (e) -- (H);
  \draw (H) -- (G) node[midway,right] {$n$};
  \draw[<->] (K) -- (e);
  \draw[<->] (E) -- (H);
  \draw[<->] (Q) -- (G);
  \draw (pp) -- (p);
\end{tikzpicture}
\end{document}
